\section{Introduction}
A stochastic response array can have small positive factorizations at
every cut without admitting one machine that realizes those
factorizations in succession. The issue is not only how many
continuations enclose the observed rows. Every chosen continuation must
also evolve by the same legal stochastic update as the other
continuations at that time. We study the resulting difference between
\emph{separate} and \emph{simultaneous} state complexity.

Our principal result gives an unbounded difference with fixed input and
output alphabets. Fix $a=1/10$, a planar rotation $R$, a seed
$x\in\{-1,1\}^2$, a word $c=(c_1,\ldots,c_N)\in\{0,1\}^N$, and a
query $j\in\{1,2\}$. The seed is read first, the commands are read in
order, and the query is supplied last. The required binary response is
\begin{equation}\label{eq:channel-intro}
 C_{N,R}(b\mid x,c,j)
       =\frac{1+b\,e_j^T R^{c_1+\cdots+c_N}a x}{2},
       \qquad b\in\{-1,1\}.
\end{equation}
The same formula is required for every input word, not merely under a
chosen distribution on words. Every entry lies between
$(1-\sqrt2/10)/2$ and $(1+\sqrt2/10)/2$.

At any one internal cut, three normalized continuations suffice and are
necessary. This remains true when the cut is required to belong to a
complete machine and all its other registers may be enlarged. Write
$W_N(R)$ for the least peak over machines satisfying all cuts at once.

\begin{theorem}[Fixed local rank and growing causal width]
\label{thm:headline}
Let $\alpha=(\sqrt5-1)/2$ and let $R$ rotate by $2\pi\alpha$.
Every internal cut of \eqref{eq:channel-intro} has separate minimum three
and ordinary normalized residual rank three. For $N\ge1$,
\begin{equation}\label{eq:headline}
 \max\left\{3,\frac16\log_2\frac{N}{24000000}\right\}
 \le W_N(R)\le
 \min\left\{3(N+1),4\lceil\sqrt{N+1}\rceil\right\}.
\end{equation}
All lower bounds allow arbitrary randomized encoders, arbitrary hidden
state continuations, and cut-dependent stochastic transition matrices.
The upper bounds realize every conditional probability exactly.

For the rational rotation
\begin{equation}\label{eq:rational-rotation}
 R_0=\frac15\begin{pmatrix}3&-4\\4&3\end{pmatrix},
\end{equation}
the same separate minima and upper bounds hold and $W_N(R_0)\to\infty$.
For each proposed state bound, an excluded horizon can be computed by
real algebraic optimization. The $3(N+1)$-state realization is rational.
\end{theorem}

The proof occupies Sections~\ref{sec:volume}--\ref{sec:quantitative}.
The logarithmic lower bound and square-root upper bound are not asserted
to be sharp. They show, with explicit constants in one family, that
uniformly bounded separate positive ranks do not control simultaneous
width even in the permissive clocked model. The rational version makes
the qualitative phenomenon independent of transcendental input data.

The principal technical point is the treatment of nonminimal hidden
representations. If a machine has $K_t$ states, its hidden state vectors
need not individually have an observable two-dimensional continuation.
Projecting those individual states into a presumed invariant polygon
would therefore restrict the competing class. Instead, let $H_t$ be the
affine span of the actually reachable state \emph{distributions}, and
intersect it with the probability simplex. Projecting
$H_t\cap\Delta_{K_t}$ gives a valid observable polytope with at most
$2^{K_t}$ vertices. Stochastic updates carry these sections forward.
A command and an idle command then force a nested sequence of observable
polytopes. A strictly positive volume increment, summed over time,
provides the lower bound. This is a necessary argument for all hidden
realizations, not a restriction to residual-state or affine-minimal
machines.

A second result concerns exact finite profile composition.
Theorem~\ref{thm:tagged-sum} proves that a verifiable retained tag makes
feasible profiles add as sets. In particular, the earlier two-cut
obstruction with frontier $(3,4),(4,3)$ yields
\[
       K_1\ge3m,\qquad K_2\ge3m,\qquad K_1+K_2\ge7m.
\]
Its optimum peak is $\lceil7m/2\rceil$, whereas either cut separately has
minimum $3m$. This closure law complements the temporal example: it
identifies an exact growing Pareto frontier, while the rotation family
keeps alphabets and all separate minima fixed. The tagged construction
has structural output zeros; it is not called a full-support example.

\subsection{Classical mechanisms and the new comparison}
Invariant polyhedral cones and normalized convex sections are classical
in positive realization \cite{Heller,BF,Vidyasagar}. Irrational rotation
as an obstruction to a finite \emph{stationary} positive realization is
also classical: see \cite[Example 4]{BF} and
\cite[Theorem 6, Example 2 and Lemma 7]{MW}. We do not claim either
mechanism as a new discovery. Reachable subspaces and observable
quotients are likewise established tools \cite[Section 4]{MW}.

The comparison proved here differs in its quantifiers. Every finite
horizon has a finite exact realization; both the state spaces and the
rows may change at every epoch. There is no common stationary matrix to
which a peripheral-spectrum obstruction can be applied. Moreover, the
least \emph{positive} factorization at each cut, not just ordinary rank,
is three. The cumulative volume bound and the explicit sequence of
finite realizations concern this finite, time-dependent problem.
Section~\ref{sec:comparison} details the relation to positive realization,
nonnegative factorization and residual automata. The explicit bounds
are claims of the present article; exhaustive priority clearance is not
inferred from the bibliography.

\begin{center}
\small
\begin{tabular}{p{.39\textwidth}p{.51\textwidth}}
\toprule
Statement & Mathematical role\\\midrule
Reachable-section and volume budget & General multi-cut lower bound; Sections 3--5\\
Rotation family & Fixed-rank, fixed-alphabet unbounded incompatibility\\
Tagged profile sum & Exact composition of feasible profile sets\\
Two-cut primitive & Inherited from revision 31; proof retained in Appendix A\\
Invariant cones and weighted replacement & Classical mechanisms; not new claims\\
\bottomrule
\end{tabular}
\end{center}
